\section{Corank budgets and relation-monoid divisors}
\label{alg:section}

The algebraic argument has two parts. First, the cap components of
idempotent matching diagrams impose a corank budget that is independent
of the length of a word in those idempotents. Second, this budget
strengthens the relation-addition construction of
\cite{OpenAILiveness2026}. We give the corner arguments and the complete
induction, so the degree bound below does not depend on importing the
source manuscript's exponential bound as an unproved assumption.

\subsection{Matching diagrams and their cap components}
\label{alg:definitions}

Write $[m]=\{1,\ldots,m\}$, with $[0]=\varnothing$.
The \emph{Brauer monoid} $\Br{m}$ consists of the perfect matchings on
two copies of $[m]$, called the left and right ports. A matching edge
between opposite sides is a \emph{transversal}; an edge whose endpoints
are on the same side is a \emph{cap}. Multiplication places the first
factor to the left of the second, identifies equally indexed ports at
the interface, and follows paths between the remaining outer ports.
Components with no outer port are discarded. Parallel edges at an
interface are retained. This is the usual diagram multiplication
associated with Brauer's algebras; see
\cite{Brauer1937,Auinger2012,EastGray2017}.

Associativity follows by contracting paths in a substring of a glued
diagram: the contraction preserves exactly the connections available
to the rest of the diagram. The identity $\Id_m$ consists of the $m$
fixed transversals. Set
\[
 \rk(a)=\#\{\text{transversals of }a\},\qquad
 \crk(a)=m-\rk(a),
\]
\[
 \supp(a)=\{i:\text{left }i\text{ is not paired with right }i\}.
\]
In particular, $\crk(a)$ is even: each side has
$\crk(a)/2$ caps. Following disjoint transversal paths through a
product gives
\begin{equation}\label{alg:rank-basic}
 \rk(ab)\le\min\{\rk(a),\rk(b)\}.
\end{equation}
The rank-$m$ diagrams are precisely permutations, and these are exactly
the units. A rank-$m$ idempotent is therefore $\Id_m$.

The \emph{cap graph} of a diagram has vertex set $[m]$: draw one red
edge for each left cap and one blue edge for each right cap, identifying
the two port labels but retaining the colors and parallel edges.
Its degrees are at most two, with at most one incident edge of each
color. A connected component is consequently an alternating path,
an alternating cycle, or an isolated vertex. A two-vertex cycle means
two parallel edges of different colors.

\begin{lemma}[Idempotent component description]\label{alg:components}
For an idempotent $e\in\Br{m}$, every cap component is either an
alternating path with an odd number of vertices, or an alternating
cycle with an even number of vertices. Each path component carries
exactly one transversal of $e$, joining its two missing-color
endpoints; an isolated vertex carries a fixed transversal. A cycle
component carries no transversal. In particular, every edge of $e$,
including every transversal, stays within one cap component.

Conversely, any collection of such colored components, completed by
these transversals, defines an idempotent. A path on $2a+1$ vertices
contributes $2a$ to corank, and a cycle on $2b$ vertices contributes
$2b$.
\end{lemma}

\begin{proof}
This is the classical component description underlying the enumeration
in \cite{Dolinka2015}; we recall its connector proof. Let $X,Y$ be the
left and right transversal endpoint sets of $e$, and let $p:X\to Y$
be the bijection specified by its transversals. At the middle interface
of $e^2$, retain the right caps of the first copy and the left caps of
the second. Contracting the edges that identify equal labels gives
the cap graph just defined. A missing blue edge is a right-transversal
endpoint in $Y$, and a missing red edge is a left-transversal endpoint
in $X$.

Because $\rk(e^2)=\rk(e)$, every endpoint of the first kind must
connect to one of the second kind. A path with both endpoints of the
same kind would join two transversals coming from the same outer
side and reduce the rank. Thus each nontrivial path has an even
number of alternating edges and an odd number of vertices. Its
endpoints define a bijection $\gamma:Y\to X$; an isolated vertex in
$X\cap Y$ is included as a connector of length zero. The transversal
bijection of $e^2$ is $p\gamma p$. Since it equals $p$, one has
$\gamma=p^{-1}$. Hence the transversal belonging to a connector path
joins that path's endpoints, and a zero-length connector forces a
fixed transversal. All remaining components are alternating cycles.
The stated component contributions now follow by counting caps.

Conversely, for a matching assembled in this way, each interface
connector is the inverse of its transversal. Gluing two copies
therefore reproduces every transversal, while the outer caps already
agree. Closed interface components do not affect the product. Thus
the square equals the original diagram.
\end{proof}

\subsection{A budget for arbitrarily long words}
\label{alg:budget-section}

\begin{theorem}[Corank budget]\label{alg:budget}
Let $e_1,\ldots,e_v\in\Br{m}$ be idempotents. Every word $w$ in
these generators, allowing repetitions and the empty word, satisfies
\begin{equation}\label{alg:budget-eq}
 \crk(w)\le\sum_{i=1}^{v}\crk(e_i).
\end{equation}
Thus each chosen generator is charged once, independently of how
often it occurs in the word.
\end{theorem}

\begin{proof}
Superimpose the cap graphs of the generators on $[m]$. Retain each
edge occurrence, including occurrences with the same endpoints from
different generators. Let $C$ be a connected component of this union
multigraph, and let $E_C$ be its number of edges, counted with these
multiplicities.

Each cap component of each individual $e_i$ is entirely contained
in a union component. By Lemma~\ref{alg:components}, it has an even
number of edges. Consequently $E_C$ is even. Connectedness gives
$E_C\ge |C|-1$, and parity improves this to
\begin{equation}\label{alg:component-budget}
 E_C\ge 2\left\lfloor\frac{|C|}{2}\right\rfloor.
\end{equation}
This includes isolated vertices, for which both sides vanish.

Every cap and every transversal of every $e_i$ stays inside a union
component, again by Lemma~\ref{alg:components}. The same is therefore
true in any glued word and in the matching obtained after contraction.
The restriction of $w$ to the two copies of $C$ is a Brauer diagram
of degree $|C|$. Its rank has the parity of $|C|$, so its corank is
at most $2\lfloor |C|/2\rfloor$, and hence at most $E_C$.
Summing over $C$ proves the result, because the total number of cap
edge occurrences is exactly $\sum_i\crk(e_i)$.
\end{proof}

\begin{corollary}[Budget relative to an idempotent]\label{alg:relative-budget}
Let $e,b_1,\ldots,b_v\in\Br{m}$ be idempotents. If $w$ is any word
in them and $z=ewe$, then
\begin{equation}\label{alg:relative-budget-eq}
 \rk(e)-\rk(z)\le\sum_{i=1}^{v}\crk(b_i).
\end{equation}
The conclusion applies to every positive power of every such
sandwich, with the same right-hand side.
\end{corollary}

\begin{proof}
Apply Theorem~\ref{alg:budget} to the generating family
$e,b_1,\ldots,b_v$, then subtract $\crk(e)$ from the resulting
ambient corank inequality. Repeated sandwiches and their powers are
still words in this fixed family.
\end{proof}

\begin{remark}
The bound does not require disjoint supports or commuting generators.
The hypothesis of idempotence supplies both facts used in the proof:
an even cap-edge count on each individual component, and preservation
of that component by transversals. Merely controlling the size of the
union of supports loses this information.
The coefficient one is sharp, already for a single nonidentity
idempotent and the word consisting of that generator. Equality with
several nonidentity generators is obtained by putting identical
left and right cap matchings on disjoint even-sized blocks and
using fixed transversals elsewhere: their product spends exactly
the sum of the block sizes.
\end{remark}

\subsection{Sharp support, equality, and stability}
\label{alg:support-section}

The same classical components give an exact refinement of the support
estimate used in \cite{OpenAILiveness2026}. This refinement is not
needed to prove Theorem~\ref{alg:budget}, but it describes precisely
where that earlier support estimate wastes rank information.

\begin{theorem}[Sharp support and its deficit]\label{alg:support}
Let $e\in\Br{m}$ be idempotent, let $c=\crk(e)$, and set
$\delta(e)=3c/2-|\supp(e)|$. Then
\begin{equation}\label{alg:support-bound}
 |\supp(e)|\le\min\{m,3c/2\}.
\end{equation}
For every admissible even $c\in[0,m]$, this is the exact maximum
support among idempotents of corank $c$.

More precisely, if the nontrivial odd path components have sizes
$2a_j+1$ with $a_j\ge1$, and the cycle components have sizes $2b_k$
with $b_k\ge1$, then
\begin{equation}\label{alg:deficit}
 \delta(e)=\sum_j(a_j-1)+\sum_k b_k.
\end{equation}
Equality $|\supp(e)|=3c/2$ holds exactly when all nontrivial
components are three-vertex paths. Removing at most $5\delta(e)$
indices, in whole components, leaves only such paths and fixed
transversals. The removed components contribute at most
$4\delta(e)$ to the original corank. The constants $5$ and $4$
are both best possible for these component deletion statements.
\end{theorem}

\begin{proof}
By Lemma~\ref{alg:components}, the support is exactly the union of
nontrivial cap components. A path of size $2a+1\ge3$ contributes
$2a$ to corank and $2a+1$ to support, giving deficit $a-1$.
A cycle of size $2b$ contributes $2b$ to both, giving deficit $b$.
Isolated vertices contribute zero. This proves
\eqref{alg:deficit}, the bound, and the equality classification.

For attainability, write $c=2q$. If $m\ge3q$, use $q$ disjoint
three-vertex path components and $m-3q$ isolated vertices. If
$2q\le m<3q$, put $r=m-2q$. Use $r$ three-vertex paths and
$q-r$ two-vertex cycles. These components use all $m$ vertices,
have rank $r$, and attain support $m$.

For stability, remove every cycle and every odd path with at least
five vertices. A removed path has $a\ge2$ and
$2a+1\le5(a-1)$; a removed cycle satisfies $2b\le5b$.
Equation~\eqref{alg:deficit} bounds the total number of removed
vertices by $5\delta(e)$. A single five-vertex alternating path
has deficit one and must be removed in its entirety under the stated
component deletion rule. Hence the factor cannot be smaller than five.
The removed paths satisfy $2a\le4(a-1)$ for $a\ge2$, while the
removed cycles satisfy $2b\le4b$. Their total corank is at most
$4\delta(e)$. The same five-vertex path has corank four and
deficit one, proving sharpness of this second constant.
\end{proof}

\begin{corollary}[Number of support maximizers]\label{alg:support-count}
For $q\ge0$ and $m\ge3q$, the number of labeled idempotents of
corank $2q$ and support $3q$ is
\begin{equation}\label{alg:maximizer-count}
 \frac{m!}{q!\,(m-3q)!}.
\end{equation}
\end{corollary}

\begin{proof}
Choose $q$ unordered disjoint triples and leave the other vertices
isolated. Their number is
\[
 \frac{m!}{(3!)^q q!\,(m-3q)!}.
\]
On each triple there are six idempotent hooks: choose the two distinct transversal
endpoints in order; the remaining vertex determines both caps.
Multiplying by $6^q$ gives the formula.
\end{proof}

\subsection{Corner reductions and unit lifts}
\label{alg:corners-section}

We next record the finite-semigroup tools needed to apply the budget
inside successive corners. A \emph{unital} homomorphism preserves
the designated identity, including when the identity of a corner
differs from that of its ambient monoid.

\begin{lemma}[Faithful reduction of a matching corner]\label{alg:corner}
Let $e\in\Br{m}$ be idempotent and put $r=\rk(e)$. There is an
injective unital homomorphism
\[
 R_e:e\Br{m}e\longrightarrow\Br{r}
\]
such that $\rk(R_e(z))=\rk(z)$ for every $z$ in the corner.
\end{lemma}

\begin{proof}
Use the endpoint sets $X,Y$ and transversal bijection $p:X\to Y$
from Lemma~\ref{alg:components}. Every $z=eae$ contains all the
left caps of $e$ on its outer left side and all the right caps of
$e$ on its outer right side. Delete those prescribed caps, retaining
the left ports in $X$, the right ports in $Y$, and every matching
edge among these retained ports. In particular, additional caps
among retained ports are not deleted.

Number $X$ by $[r]$ and give the right port $p(x)$ the same number
as the left port $x$. The retained diagram defines $R_e(z)$.
Restoring the prescribed caps reconstructs $z$, proving injectivity;
no transversal was deleted, proving rank preservation.

At the interface of $zw$ with $z,w\in e\Br{m}e$, the prescribed
caps form precisely the connector of $e^2$. The connector sends
$p(x)$ to $x$, as established in Lemma~\ref{alg:components}.
Contracting it therefore identifies equal numbers in the two reduced
diagrams. Any closed connector components are irrelevant. This gives
$R_e(zw)=R_e(z)R_e(w)$, and the numbering gives $R_e(e)=\Id_r$.
If $r=0$, all outer ports are prescribed caps, so the same argument
identifies the corner with the one-element monoid $\Br{0}$.
\end{proof}

This explicit construction is a version of the local-monoid methods
used for diagram semigroups; compare
\cite{Auinger2012,Pin1997,OpenAILiveness2026}.

\begin{lemma}[Minimum-rank corners lift units]\label{alg:unit-lifts}
Let $S\subseteq\Br{m}$ be nonempty and closed under multiplication,
and let $\phi:S\twoheadrightarrow M$ be multiplicative, where $M$
is a monoid. For each idempotent $q\in M$ there is an idempotent
$e\in S$ above $q$. Every such $e$ gives a unital surjection
\begin{equation}\label{alg:induced-corner-map}
 R_e(eSe)\twoheadrightarrow qMq,\qquad R_e(z)\longmapsto\phi(z).
\end{equation}
If $e$ has minimum rank among idempotents above $q$, every element
of $eSe$ whose image is a unit of $qMq$ reduces to a permutation
diagram. In particular all units of $qMq$ have unit lifts in the
reduced corner.
\end{lemma}

\begin{proof}
In any finite semigroup an element has a positive idempotent power:
its powers are eventually periodic, and an exponent sufficiently
large and divisible by the eventual period has this property.
An idempotent power of a preimage of $q$ is still above $q$.
Furthermore, $\phi(eSe)=qMq$ by surjectivity, and the identity
$e$ of the corner maps to $q$. Lemma~\ref{alg:corner} transfers this
map faithfully to the reduced corner, proving
\eqref{alg:induced-corner-map}.

Now let $y\in eSe$ map to a unit. Some positive power of its image
is $q$, and a further positive power of $y$ is an idempotent $z$
above $q$. Minimality and \eqref{alg:rank-basic} give
\[
 \rk(e)\le\rk(z)\le\rk(y)\le\rk(e).
\]
Thus $R_e(y)$ has full rank and is a permutation. Its inverse is
a positive power of itself and therefore belongs to $R_e(eSe)$.
Surjectivity supplies such a preimage for every target unit.
\end{proof}

For a finite set $H$, let $\Rel{H}$ be the monoid of all binary
relations on $H$, in path order:
$(x,z)\in AB$ when $(x,y)\in A$ and $(y,z)\in B$ for some $y$.
Its identity is $\Id_H$ and its zero is the empty relation. Its
units are exactly the permutation relations. A relation on
$F\subseteq H$ is identified with the same subset of $H\times H$;
equivalently, it is a relation $A$ satisfying $A=\Id_FA\Id_F$.

\begin{lemma}[A full relation corner]\label{alg:relation-corner}
Suppose $P\in\Rel{H}$ is idempotent and $F\subseteq H$ satisfies
$\Id_FP\Id_F=\Id_F$. Put $E=\Id_F$ and $Q=PEP$. Then $Q$ is
idempotent, $PQ=QP=Q$, and
\begin{equation}\label{alg:relation-corner-iso}
 Q\Rel{H}Q\longrightarrow\Rel{F},\qquad Z\longmapsto EZE
\end{equation}
is a unital isomorphism, with inverse $A\mapsto PAP$.
\end{lemma}

\begin{proof}
From $P^2=P$ and $EPE=E$ one obtains
$Q^2=Q$, $PQ=QP=Q$, $EQ=EP$, $QE=PE$, and $EQE=E$.
If $Z=QZQ$, then $PZP=Z$ and
\[
 Z=PEPZPEP=PEZEP.
\]
Thus $EZE$ determines $Z$. Conversely, for $A=EAE$,
\[
 QAQ=PAP,\qquad E(PAP)E=A.
\]
These formulas prove the inverse assertion. For $A,B\in\Rel{F}$,
\[
 (PAP)(PBP)=PAPBP=PA(EPE)BP=PABP,
\]
so the bijection respects products. It sends the corner identity
$Q$ to $E$, proving unitality.
\end{proof}

\subsection{Amplifying the rank cost of one relation addition}
\label{alg:amplification-section}

The next construction follows the two-stage relation-corner strategy
of \cite{OpenAILiveness2026}. Its rank accounting is different:
Corollary~\ref{alg:relative-budget} controls both the complete chain
and each local sandwich, without passing through support sizes.

\begin{definition}\label{alg:admissible-bound}
A number $L\ge0$ is an \emph{addition bound at size $h\ge2$} if the
following holds. Whenever an identity-containing submonoid
$S\subseteq\Br{m}$ admits a unital surjection
$\phi:S\twoheadrightarrow\Rel{H}$, with $|H|=h$ and with permutation
diagram lifts for all permutations of $H$, every idempotent $a\in S$
with
\[
 \phi(a)=\Id_H\cup\{(x,y)\},\qquad x\ne y,
\]
has $\crk(a)\ge L$.
\end{definition}

\begin{proposition}[General amplification]\label{alg:amplification}
Let $p,q\ge1$, $h\ge p+q+1$, and $k=h-p-q+1$. If $L$ is an
addition bound at size $k$, then
\begin{equation}\label{alg:general-amplification}
 \frac{pq}{p+q}\,L
\end{equation}
is an addition bound at size $h$.
\end{proposition}

\begin{proof}
Fix data from Definition~\ref{alg:admissible-bound} at size $h$ and
write $c=\crk(a)$. Choose disjoint sets
$U=\{u_1,\ldots,u_p\}$, $V=\{v_1,\ldots,v_q\}$, and $\{z\}$
inside $H$. Conjugating $a$ by permutation lifts supplies idempotents
$s_i,t_j\in S$ such that
\begin{equation}\label{alg:conjugate-images}
 \phi(s_i)=\Id_H\cup\{(u_i,z)\},\qquad
 \phi(t_j)=\Id_H\cup\{(z,v_j)\}.
\end{equation}
All these conjugates have corank $c$. The inverses of the chosen
permutation lifts belong to $S$, since each is a power of a
finite-order permutation.

Put $P_0=\Id_{H\setminus\{z\}}$ and choose an idempotent lift
$b_0\in S$ of $P_0$ by taking a positive idempotent power.
For each pair $(i,j)$ let $g_{ij}=b_0s_it_jb_0$. The path through
$z$ in \eqref{alg:conjugate-images}, followed by deletion of its
endpoints from the outer relation, gives
\[
 \phi(g_{ij})=P_0\cup\{(u_i,v_j)\}.
\]
List the $pq$ pairs in any order. If $E_\ell\subseteq U\times V$
contains the first $\ell$ listed pairs, define $P_\ell=P_0\cup E_\ell$.
Because $U\cap V=\varnothing$,
\begin{equation}\label{alg:union-relations}
 (P_0\cup A)(P_0\cup B)=P_0\cup(A\cup B)
 \quad(A,B\subseteq U\times V).
\end{equation}
In particular all $P_\ell$ are idempotent.

Inductively choose $b_\ell$ to be a positive idempotent power of
$b_{\ell-1}g_{ij}b_{\ell-1}$, where $(i,j)$ is the $\ell$th pair.
Equation~\eqref{alg:union-relations} gives
\begin{equation}\label{alg:chain}
 b_\ell^2=b_\ell,\qquad
 b_\ell\in b_{\ell-1}Sb_{\ell-1},\qquad
 \phi(b_\ell)=P_\ell.
\end{equation}
Write $r_\ell=\rk(b_\ell)$ and
$\Delta_\ell=r_{\ell-1}-r_\ell\ge0$. Every $b_\ell$ is a word in
the fixed idempotent family $b_0,s_1,\ldots,s_p,t_1,\ldots,t_q$
and lies in the corner of $b_0$. Corollary~\ref{alg:relative-budget}
therefore gives the single global budget
\begin{equation}\label{alg:global-budget}
 \sum_{\ell=1}^{pq}\Delta_\ell
 =r_0-r_{pq}\le(p+q)c.
\end{equation}

We prove that each $\Delta_\ell\ge L$. Fix the pair $(i,j)$ at
step $\ell$, put $b=b_{\ell-1}$ and $P=P_{\ell-1}$, and pass to
the reduced corner
\[
 T=R_b(bSb)\subseteq\Br{r_{\ell-1}},\qquad
 \beta=R_b(b_\ell).
\]
The induced map $\theta:T\twoheadrightarrow P\Rel{H}P$ is unital
by Lemma~\ref{alg:unit-lifts}, and $\beta$ is idempotent with
\[
 \theta(\beta)=P_\ell,\qquad \crk(\beta)=\Delta_\ell.
\]
Here and below corank is computed in the currently displayed ambient
Brauer monoid; all ranks are preserved by the reductions.

Retain precisely
\[
 F=\{u_i,v_j\}\cup
    \bigl(H\setminus(U\cup V\cup\{z\})\bigr),\qquad E=\Id_F.
\]
Then $|F|=k$. The only potential earlier added pair with both
endpoints in $F$ is $(u_i,v_j)$, and this pair has not yet been added.
Thus $EPE=E$. Put $Q=PEP$. By Lemma~\ref{alg:relation-corner},
restriction identifies $Q\Rel{H}Q$ with $\Rel{F}$, and
\[
 Q(P\Rel{H}P)Q=Q\Rel{H}Q.
\]
Moreover, $PP_\ell=P_\ell P=P_\ell$, so
\begin{equation}\label{alg:new-edge-survives}
 E(QP_\ell Q)E=EP_\ell E
 =\Id_F\cup\{(u_i,v_j)\}.
\end{equation}

Choose an idempotent $d\in T$ of minimum rank among those with
$\theta(d)=Q$, and set $r'=\rk(d)$. Lemma~\ref{alg:unit-lifts}
and the restriction isomorphism give a unital surjection
\[
 \psi:T'=R_d(dTd)\twoheadrightarrow\Rel{F}
\]
with permutation diagram lifts for every permutation of $F$.
Take a positive idempotent power $v$ of $R_d(d\beta d)$.
Equation~\eqref{alg:new-edge-survives} says that its image before
taking the power is already the idempotent relation
$\Id_F\cup\{(u_i,v_j)\}$; the power has the same image.
The size-$k$ addition bound consequently gives $\crk(v)\ge L$.

On the other hand, apply Corollary~\ref{alg:relative-budget} to the
two idempotents $d,\beta$ in $\Br{r_{\ell-1}}$. The corresponding
power of $d\beta d$ loses at most $\crk(\beta)$ relative to
$\rk(d)$. Faithfulness and rank preservation of $R_d$ give
\[
 L\le\crk(v)\le\crk(\beta)=\Delta_\ell.
\]
This proves the local claim. Summing it and using
\eqref{alg:global-budget} yields $pqL\le(p+q)c$, as required.
\end{proof}

\begin{theorem}[Exponential cost of an addition]\label{alg:rank}
For every $h\ge2$, the number
\begin{equation}\label{alg:F}
 F(h)=2\left(\frac52\right)^{\lfloor(h-2)/9\rfloor}
\end{equation}
is an addition bound at size $h$.
\end{theorem}

\begin{proof}
Any idempotent above a nonidentity relation has positive corank:
otherwise it is the ambient identity, contradicting unitality.
Corank is even, so it is at least two. This proves the assertion
for $2\le h\le10$. For $h\ge11$, use strong induction and
Proposition~\ref{alg:amplification} with $p=q=5$. It gives an
addition bound $(5/2)F(h-9)=F(h)$. All smaller sets used in this
step have size $h-9\ge2$.
\end{proof}

\begin{theorem}[Degree of a full relation-monoid divisor]\label{alg:divisor}
Let $H$ have $h\ge2$ elements. If an identity-containing submonoid
$S\subseteq\Br{k}$ has a unital surjection onto $\Rel{H}$, then
\begin{equation}\label{alg:divisor-bound}
 k\ge F(h)=2\left(\frac52\right)^{\lfloor(h-2)/9\rfloor}.
\end{equation}
No assumption about permutation lifts is needed.
\end{theorem}

\begin{proof}
Choose an idempotent $e\in S$ of minimum rank above $\Id_H$;
the identity of $S$ is one possible choice. Its reduced corner
$R_e(eSe)\subseteq\Br{r}$, where $r=\rk(e)$, maps unitally onto
$\Rel{H}$ and has permutation lifts by Lemma~\ref{alg:unit-lifts}.
Lift any relation $\Id_H\cup\{(x,y)\}$ with $x\ne y$ to that
corner and take an idempotent power. Theorem~\ref{alg:rank} bounds
its corank below by $F(h)$, whereas its corank is at most $r\le k$.
\end{proof}

\subsection{Choice of parameters and finite integer bounds}
\label{alg:optimization-section}

\begin{proposition}[Optimal fixed rectangular amplification]
\label{alg:optimal-parameters}
Among fixed positive integer choices of $p,q$, the ratio
\begin{equation}\label{alg:rate}
 \frac{\log\bigl(pq/(p+q)\bigr)}{p+q-1}
\end{equation}
is maximized at $p=q=5$. Thus the fixed choice in
Theorem~\ref{alg:rank} has the best asymptotic exponential rate
within Proposition~\ref{alg:amplification}.
\end{proposition}

\begin{proof}
At fixed $p+q$, balancing $p,q$ maximizes their product. For an even
sum write $p=q=t$; the rate is
$f(t)=\log(t/2)/(2t-1)$. Its derivative has the sign of
\[
 N(t)=2-1/t-2\log(t/2),\qquad
 N'(t)=(1-2t)/t^2<0\quad(t>1/2).
\]
One has $N(4)>0>N(5)$: for example, $\log 2<3/4$ and
$\log(5/2)>312/343>9/10$, using the first two positive terms of
$\log((1+x)/(1-x))$ at $x=3/7$ for the latter bound.
Thus only $t=4,5$ can maximize the rate
over integers. The inequality $f(5)>f(4)$ is equivalent to
$5^7>2^{16}$, namely $78125>65536$.

For an odd sum $2t+1$, the best factor is $t(t+1)/(2t+1)$, and
\[
 \frac{t(t+1)}{2t+1}
 \le\sqrt{\frac t2\frac{t+1}{2}}.
\]
Its deletion cost $2t$ is the arithmetic mean of $2t-1$ and
$2t+1$. Taking logarithms shows that its rate is at most the
larger of the rates for the two neighboring even sums. This
completes the optimization.
\end{proof}

The resulting base-two rate is
\[
 \frac{\log_2(5/2)}9=0.1468808994\ldots.
\]
The optimization concerns this particular rectangular construction;
it is not an optimality assertion about the true minimum degree of
a Brauer monoid admitting the required divisor.

\begin{proposition}[Parity-aware finite recurrence]\label{alg:finite-recurrence}
Define positive integers $D(h)$ recursively by $D(2)=1$ and, for
$h\ge3$,
\begin{equation}\label{alg:D}
 D(h)=\max\left\{1,\ \max_{2\le s\le h-1}
 \left\lceil\frac{\lfloor s^2/4\rfloor}{s}\,
 D(h-s+1)\right\rceil\right\}.
\end{equation}
Then $2D(h)$ is an addition bound at size $h$, and the conclusion
of Theorem~\ref{alg:divisor} can be strengthened to $k\ge2D(h)$.
\end{proposition}

\begin{proof}
Use strong induction. For each allowed $s$, choose balanced positive
$p,q$ with $p+q=s$ and $pq=\lfloor s^2/4\rfloor$.
Proposition~\ref{alg:amplification} and the induction hypothesis
give corank at least $2pqD(h-s+1)/s$. Corank is an even integer,
so it is at least twice the corresponding ceiling in
\eqref{alg:D}. The baseline bound is two, and taking the maximum
proves the addition claim. The minimum-rank corner proof of
Theorem~\ref{alg:divisor} then applies unchanged.
Every recursive argument $h-s+1$ lies between $2$ and $h-1$;
no monotonicity of $D$ is assumed or needed.
\end{proof}
